\section{Network Forensics}

Network forensics involves the capture, recording, and analysis of network traffic to investigate security incidents or legal evidence.\\
Unlike static file system forensics, network forensics deals with volatile data and the consequences of actions performed across digital networks.

\subsection{OSINT Integration}
Modern investigations adopt a multi approach, combining traditional packet analysis with \textbf{Open Source Intelligence (OSINT)}.\\
OSINT serves as a supplement to digital forensic data by gathering information from publicly available sources to reconstruct a suspect's digital fingerprint.

\subsubsection*{Key OSINT Evidence Sources}
\begin{itemize}
    \item \textbf{Archived Snapshots:} Using tools like the \textit{Wayback Machine} to access historical versions of websites that are no longer online.
    \item \textbf{Domain Registration:} Checking the age and ownership of domains. Recently registered domains can often be a "red flag" for malicious campaigns.
    \item \textbf{Public Repositories and Forums:} Analyzing code leaks, forum posts, and usernames to identify affiliations or signs of compromise.
    \item \textbf{Social Media Profiling:} Tracking activity patterns, location data from posts, and sentiment analysis to build a psychological or behavioral profile of a suspect.
\end{itemize}

\subsection{Technical Acquisition and Analysis}
The primary objective in network forensics is the reconstruction of communication timelines. This process is often complicated by encryption (TLS, SSH, IPsec) and obfuscation.

\subsubsection*{Analysis Workflow}
When deep packet inspection (DPI) is under encryption, investigators must use to complementary techniques:
\begin{enumerate}
    \item \textbf{Lower-Layer Analysis:} Accessing metadata at the data link, network, and transport layers to identify IP addresses and connection timestamps.
    \item \textbf{Endpoint Integration:} Accessing the host device to recover symmetric session keys, allowing for the decryption of TLS packets with tools like Wireshark.
    \item \textbf{Pattern Recognition:} Identifying specific traffic behaviors even in the absence of clear-text content.
\end{enumerate}

\subsubsection{Tool for Network Investigation}
\begin{table}[H]
\centering
\small
\begin{tabular}{|l|p{8.5cm}|}
\hline
\textbf{Tool} & \textbf{Forensic Application} \\
\hline
\textbf{Wireshark} & Industry standard for packet analysis and decryption using recovered session keys. \\
\hline
\textbf{Nmap} & Used to identify active services or "shadow servers" on a node and corroborate pcap data. \\
\hline
\textbf{Xplico / NetworkMiner} & Specialized in "carving" application-level data (e.g., extracting files, web pages, or emails) from network streams. \\
\hline
\textbf{PCAP Format} & The standard data format for storing captured packets, preserving the raw evidence for timeline reconstruction. \\
\hline
\end{tabular}
\caption{Primary tools and formats in Network Forensics.}
\end{table}

\subsection{Network Forensic Phases}

The lifecycle of a network forensic investigation ensures that volatile network data is captured, analyzed for patterns, and documented in a manner that remains defensible in legal contexts.

\subsubsection{Data Preservation and Snapshots}
Because network information is often outside the investigator's direct control (residing on ISP or third-party servers), data can be modified, moved, or deleted at any time.

\subsubsection*{Methodology for Admissibility}
To preserve the "original state" of network evidence for legal use, investigators must adopt snapshot-driven tools that:
\begin{itemize}
    \item \textbf{Capture Snapshots:} Generate immediate captures of web archives, social media posts, and public repositories.
    \item \textbf{Digital Signing:} Apply cryptographic signatures and timestamps to verify that the information existed in that specific state at that specific time.
    \item \textbf{Metadata Retrieval:} Replicate associated metadata (IP headers, server logs) to provide a verifiable chain of evidence.
\end{itemize}

\subsubsection{Evidence Analysis and Relationship Mapping}
The analysis phase focuses on identifying relationships between disparate digital artifacts and reconstructing the attack timeline.

\subsubsection*{Correlation Graphs}
Analysis typically involves the creation of a graph model to identify connections between:
\begin{itemize}
    \item \textbf{Entities:} Relationships between individuals, usernames, and physical locations.
    \item \textbf{Network Artifacts:} Correlation of IP addresses, domain names, and email metadata.
    \item \textbf{Attack Vectors:} Tracing the origin of a digital incident (e.g., distinguishing between a legitimate email and a forged bot-generated spam campaign).
\end{itemize}

\noindent
This process often requires cooperation with third parties, such as network administrators or service providers, to access restricted metadata and identify patterns of malicious behavior.

\subsubsection{Example: IP Geolocation (OSINT)}
A common forensic task is trying to determine the geographical location of a node or user using public databases.

\subsubsection*{Contextual IP Tools}
Tools such as \textbf{IPinfo}, \textbf{ip-api.com}, and \textbf{MaxMind GeoIP} allow investigators to immediately retrieve:
\begin{itemize}
    \item \textbf{Geographical Data:} City, region, and approximate GPS coordinates.
    \item \textbf{Organization Data:} ASN (Autonomous System Number) and the hosting provider or ISP.
    \item \textbf{Obfuscation:} Identifying if an IP is associated with a VPN, proxy, or Tor exit node.
\end{itemize}


\subsubsection{Reconstruction and Timeline Correlation}
Effective network forensics requires the synchronization of multiple data.

\subsubsection*{Pillars of Reconstruction}
\begin{itemize}
    \item \textbf{Causal Correlation:} Establishing a direct link between an action  and its result.
    \item \textbf{Temporal Synchronization:} Aligning timestamps from varied sources (server logs, ISP metadata, local machine clocks). 
    \item \textbf{Identity Mapping:} Linking network identifiers (IP addresses, MAC addresses, session tokens) to physical entities or specific user accounts.
\end{itemize}

\noindent
Evidence retrieved from the network is categorized based on its source and nature, influencing how it is handled in court.

\begin{table}[H]
\centering
\small
\begin{tabular}{|p{3.5cm}|p{8.5cm}|}
\hline
\textbf{Evidence Type} & \textbf{Investigative Value} \\
\hline
\textbf{Full Content Data} & Complete PCAP captures. Provides total visibility but is storage-intensive and often limited by encryption. \\
\hline
\textbf{Flow Data (NetFlow)} & Metadata summary of connections (source, destination, duration, bytes). Crucial for identifying traffic patterns and anomalies without looking at content. \\
\hline
\textbf{Transactional Data} & Interaction logs from specific services (HTTP logs, DNS queries). Useful for identifying the "intent" of a user session. \\
\hline
\end{tabular}
\caption{Taxonomy of Network Evidence.}
\end{table}

\subsubsection{Challenges in Modern Network Analysis}
As networks evolve, forensics must adapt to increasing layers of complexity.

\subsubsection*{The Encryption and Privacy Paradox}
The widespread adoption of end-to-end encryption and VPNs has shifted the focus of network forensics from \textbf{content analysis} to \textbf{traffic analysis}.


\subsubsection{Final Reporting}
The ultimate goal of the investigation is to produce a document that answers the fundamental forensic questions: \textbf{Who, What, Where, When, and Why}.

\subsubsection*{Operational Reporting Guidelines}
\begin{itemize}
    \item \textbf{Objectivity:} The report must strictly adhere to the facts observed in the data, avoiding speculative leaps.
    \item \textbf{Reproducibility:} Another investigator should be able to follow the steps outlined in the report and reach the same conclusions using the same raw evidence.
    \item \textbf{Executive Summary:} A high-level overview that translates technical PCAP data into a narrative understandable by non-technical stakeholders (judges, lawyers, managers).
\end{itemize}

\subsection{Social Media Forensics}

Social Media Forensics involves the identification, preservation, and analysis of data generated on social networking platforms.\\
Given the dynamic nature of these environments, evidence is highly volatile and requires specialized acquisition techniques to maintain its integrity for judicial purposes.

\subsubsection{Data Acquisition and Preservation}
Unlike static local storage, social media data is hosted on third-party infrastructures, making traditional bit-stream imaging impossible. Investigators must rely on a combination of API-based collection and advanced scraping techniques.



\subsubsection*{Methodologies for Legal Admissibility}
\begin{itemize}
    \item \textbf{Remote Snapshot Capture:} Using forensic tools to take authenticated "snapshots" of profiles, posts, and comments. These must include cryptographic hashing and server-side timestamps.
    \item \textbf{Metadata Extraction:} Capturing hidden details such as the \textit{User Agent}, original upload timestamps, and associated IP addresses where available.
    \item \textbf{Chain of Custody for Web Evidence:} Documenting the entire path of acquisition, including the URL, the tools used, and the digital signatures applied at the moment of capture.
\end{itemize}

\subsubsection{Forensic Artifacts in Social Media}
The analysis focuses on reconstructing user behavior and identifying the origin of digital content.

\begin{table}[H]
\centering
\small
\begin{tabular}{|p{3.5cm}|p{8cm}|}
\hline
\textbf{Artifact Category} & \textbf{Investigative Relevance} \\
\hline
\textbf{Profile Information} & Verified identities, bio changes, and linked accounts (cross-platform OSINT). \\
\hline
\textbf{Interaction Logs} & Timeline of likes, shares, and comments to establish intent or involvement in specific campaigns. \\
\hline
\textbf{Geotagging and EXIF} & Analysis of location data embedded in posts or images to verify the physical presence of a suspect. \\
\hline
\textbf{Direct Messages (DM)} & Private communication logs that often require specific legal warrants or the recovery of session tokens. \\
\hline
\end{tabular}
\caption{Classification of key social media forensic artifacts.}
\end{table}

\subsubsection{Challenges and Countermeasures}
Modern platforms implement several features that complicate forensic investigations, requiring a pivot toward behavioral and pattern-based analysis.

\subsubsection*{Anti-Forensic Obstacles}
\begin{itemize}
    \item \textbf{Ephemeral Content:} Posts that automatically disappear (e.g., IG Stories). 
    \item \textbf{Encryption and Privacy Settings:} End-to-end encryption in messaging forcing investigators to focus on physical endpoint acquisition (smartphone).
    \item \textbf{Account Deletion and Spoofing:} The ease of creating "sockpuppet" accounts necessitates the use of OSINT to correlate IP addresses and behavioral fingerprints across different sessions.
\end{itemize}

\noindent


\subsection{Threat Intelligence and Malware Forensics}
In the context of modern cyber-attacks, identifying a single instance of malware is often insufficient. Digital forensics must integrate \textbf{Threat Intelligence (TI)} to understand the broader scope of an intrusion and attribute it to specific threat actors or campaigns.

\subsubsection*{Indicators of Compromise (IoC)}
IoCs are technical artifacts that provide evidence of a potential security breach. In forensic analysis, these indicators are crucial for cross-referencing local findings with global intelligence.

\begin{table}[H]
\centering
\small
\begin{tabular}{|p{3cm}|p{8.5cm}|}
\hline
\textbf{IoC Category} & \textbf{Examples and Forensic Utility} \\
\hline
\textbf{Host-Based} & File hashes (MD5, SHA-256), specific file paths, and registry keys associated with persistence mechanisms. \\
\hline
\textbf{Network-Based} & Malicious URLs, C2 (Command and Control) domain names, and suspicious IP addresses. \\
\hline
\textbf{Behavioral} & Patterns of activity, such as unusual outbound traffic or unauthorized use of administrative tools. \\
\hline
\end{tabular}
\caption{Classification of Indicators of Compromise.}
\end{table}

\subsubsection*{Malware Families and Campaign Attribution}
Most malware is not unique but belongs to established "families" that share architectural traits. Forensic analysis of malware operations focuses on several key objectives:
\begin{enumerate}
    \item \textbf{Functionality Analysis:} Determining how the malware operates in the wild and identifying its intended impact (e.g., data exfiltration, ransomware).
    \item \textbf{Code Similarity and Reuse:} Identifying reused modules or techniques to attribute the malware to a specific Advanced Persistent Threat (APT) group.
    \item \textbf{Infrastructure Mapping:} Gathering information on the Command and Control (C2) infrastructure to understand how the malware communicates and receives updates.
\end{enumerate}

\begin{figure}[H]
\centering
\begin{tikzpicture}[font=\small,node distance=6mm,>=Latex,
box/.style={rectangle,rounded corners,draw=black,fill=gray!5,align=center,text width=4cm,minimum height=8mm}]
\node[box] (collect) {\textbf{Intelligence Gathering}\\OSINT, Virustotal, AlienVault};
\node[box, below=of collect] (anal) {\textbf{Forensic Analysis}\\Behavioral \& Static};
\node[box, below left=8mm and -5mm of anal] (ioc) {\textbf{IoC Extraction}\\Hashes, IPs, Domains};
\node[box, below right=8mm and -5mm of anal] (attr) {\textbf{Attribution}\\Campaigns \& Families};
\draw[->] (collect) -- (anal);
\draw[->] (anal) -- (ioc);
\draw[->] (anal) -- (attr);
\end{tikzpicture}
\caption{The relationship between forensic analysis and threat intelligence integration.}
\end{figure}

\noindent
Platforms like \textbf{VirusTotal} and \textbf{AlienVault OTX} allow for a collaborative approach to forensics. By sharing IoCs and malware hashes, organizations can identify if a local infection is part of a global campaign, enabling a more rapid and informed defensive response. \\
This information is crucial for tracking \textit{metamorphisis} of malwares that alters its code to evade signature-based detection.

\subsubsection{Focus: Tor Network}

Tor is an open-source anonymity network designed to protect users against traffic analysis and network surveillance. It operates by routing network traffic through a distributed, multi-layered volunteer relay network, obscuring both the origin and the destination of the data streams.

\subsubsection*{Architectural Principles and Directory Authorities}
The establishment of a Tor circuit relies on an asymmetrical trust model overseen by a dedicated management infrastructure.
\begin{itemize}
    \item \textbf{Directory Authorities (DA):} A fixed set of trusted redundant servers (typically 9 or 10 globally) responsible for maintaining and publishing the network consensus. They continuously monitor relay availability, bandwidth, and security flags.
    \item \textbf{Network Consensus:} A signed descriptor file published hourly by the DAs. Before establishing a connection, the Tor client fetches this consensus to obtain an authenticated, up-to-date list of all active and reliable relays within the network.
\end{itemize}

\subsubsection*{Three Layer Relay Mechanism}
A standard Tor circuit consists of three distinct nodes randomly selected by the client from the network consensus. This multi-hop structure ensures that no single node can link the user's IP address to the destination server.

\begin{itemize}
    \item \textbf{Entry / Guard Node:} Sees the client's real IP address,cannot see the destination server.\\
    It can only decrypts the outermost encryption layer to identify the Middle Node.

    \item \textbf{Middle Node} Sees only the Entry Node and the Exit Node.\\
    It decrypts the second layer of encryption and has zero visibility of endpoints.

    \item \textbf{Exit Node:} Sees the Middle Node and the unencrypted destination server target. 
\end{itemize}



\subsubsection*{Multi-Layered Onion Encryption}
The core cryptographic mechanism relies on symmetric keys established via a DH handshake.
\begin{itemize}
    \item \textbf{Layered Encapsulation:} The client encrypts the original payload multiple times, creating nested layers corresponding to each relay in the circuit.
    \item \textbf{Hop-by-Hop Decryption:} As a packet moves through the network, each node removes its respective layer of encryption using its symmetric key, revealing only the routing instructions for the next hop.
    \item \textbf{Exit Node Vulnerabilities:} Because the Exit Node strips away the final layer of Tor encryption, it can observe raw application data (e.g., HTTP traffic). If upper-layer protocols lack end-to-end encryption (such as HTTPS/TLS), the Exit Node operator can capture sensitive credentials or alter payloads.
\end{itemize}

\begin{figure}[H]
\centering
\begin{tikzpicture}[font=\small,node distance=15mm,>=Latex,
box/.style={rectangle,rounded corners,draw=black,fill=gray!5,align=center,text width=2.8cm,minimum height=8mm}]
\node[box] (client) {\textbf{Tor Client}\\Fetches Consensus};
\node[box, right=of client] (guard) {\textbf{Entry / Guard}\\Sees Client IP};
\node[box, right=of guard] (mid) {\textbf{Middle Node}\\Isolated Metadata};
\node[box, right=of mid] (exit) {\textbf{Exit Node}\\Sees Destination};
\node[box, below=of exit, yshift=-1cm] (dest) {\textbf{Destination Server}\\(Clear Web)};
\draw[->, thick] (client) -- (guard) node[midway, above]{Layer 3};
\draw[->, thick] (guard) -- (mid) node[midway, above]{Layer 2};
\draw[->, thick] (mid) -- (exit) node[midway, above]{Layer 1};
\draw[->, thick, red] (exit) -- (dest) node[midway, left]{Unencrypted};
\end{tikzpicture}
\caption{Cryptographic data transit and node visibility within a Tor circuit.}
\end{figure}

\subsubsection*{Tor Bridges and Obfuscation Mechanisms}
To counter censorship and network blocking (e.g., national firewalls blocking public Tor Entry node IP listings), the network utilizes non-public entry points.
\begin{itemize}
    \item \textbf{Tor Bridges:} Hidden relays that are not published in the standard Directory Authority consensus. Users must request bridge addresses out-of-band via secure channels.
    \item \textbf{Pluggable Transports (obfs4, Meek):} Obfuscation protocols that alter Tor traffic patterns. They transform the cryptographic packet signature into seemingly randomized bytes or encapsulate it inside legitimate protocols (like standard HTTPS or cloud service traffic) to bypass Deep Packet Inspection (DPI) blocks.
\end{itemize}

\subsubsection*{Onion Services and Hidden Volumes}
Tor enables end-to-end anonymity through \textbf{Onion Services} (formerly known as Hidden Services, identified by \texttt{.onion} domains). Traffic never leaves the Tor network, eliminating the security risks associated with public Exit Nodes.
\begin{itemize}
    \item \textbf{Rendezvous Points:} To connect a client to an Onion Service without exposing either party's IP address, the connection is negotiated anonymously at a mutually agreed-upon Middle Node known as a Rendezvous Point.
    \item \textbf{Forensic Footprint:} From a forensic standpoint, tracing a connection to an active onion service is mathematically unfeasible via standard network monitoring. Investigations must pivot toward host-based artifacts or application-layer vulnerabilities.
\end{itemize}

\subsubsection*{Forensic Artifacts and Investigative Strategies}
Despite Tor's robust transport layer anonymity, investigators can collect significant circumstantial evidence from host machines and boundary infrastructure.

\begin{itemize}
    \item \textbf{Volatile Memory Artifacts:} The Tor Browser utilizes a customized version of Firefox engineered to minimize local storage footprints. However, plaintext session details, visited URLs, and cryptographic configurations remain stored within the client machine's volatile memory (RAM) until a power cycle occurs.
    \item \textbf{Local File System Artifacts:} While persistent caching is disabled by default, unexpected system behavior (e.g., Windows pagefile/swap file writes, system crashes, or unallocated space residuals) may contain strings or fragments associated with Tor execution.
    \item \textbf{Network-Edge Artifacts:} Local ISPs can easily detect that a user is connected to the Tor network by cross-referencing public Entry Node IP addresses found in the Directory Authority consensus, unless advanced pluggable transports are deployed.
\end{itemize}
